%! The Rational Parent Function \& Transformations — Unit 4, Lesson 4.4 — Warm-Up
\documentclass[10pt]{article}
\usepackage{algebra2-article}
\usepackage{algebra2-boxes}
\newcommand{\TallMath}[1]{$\displaystyle #1\rule[-1.4em]{0pt}{3.2em}$}

\begin{document}

\pageheader{Unit 4, Lesson 4.4}{Warm-Up}
\namedateperiod

\vspace{0.10in}

\begin{notesbox}{Getting Ready: Skills We'll Use Today}
\small Three skills, one per leg of today's lesson. Work quickly.

\vspace{4pt}
\textbf{1. The reciprocal, up close and far out} (Lesson 4.0). Complete the table for
$y=\dfrac{1}{x}$.
\begin{center}
\renewcommand{\arraystretch}{1.35}
\begin{tabular}{|c|c|c|c|c|c|c|c|}
\hline
$x$ & $-4$ & $-1$ & $-\frac12$ & $0$ & $\frac12$ & $1$ & $4$ \\ \hline
$y=\frac1x$ & \blank{1.0cm} & \blank{1.0cm} & \blank{1.0cm} & \blank{1.0cm} & \blank{1.0cm} & \blank{1.0cm} & \blank{1.0cm} \\ \hline
\end{tabular}
\end{center}
\begin{enumerate}[label=(\alph*), leftmargin=2.1em, itemsep=3pt]
  \item Domain: \blank{3.0cm} \quad Range: \blank{3.0cm}
  \item As $x$ gets close to $0$ the outputs \blank{2.6cm}, so the graph has a \emph{vertical
        asymptote} with equation \blank{1.4cm}.
  \item As $x$ runs far out to the right or left the outputs \blank{2.6cm}, so the graph has a
        \emph{horizontal asymptote} with equation \blank{1.4cm}.
\end{enumerate}

\vspace{4pt}
\textbf{2. Transformations you already own} (Units 1--2). Name the turning point of each.
\begin{enumerate}[label=(\alph*), leftmargin=2.1em, itemsep=3pt]
  \item $y=|x-3|+1$: \; vertex $(\blank{0.7cm},\,\blank{0.7cm})$ --- the parent moved
        \blank{1.6cm} $3$ and \blank{1.4cm} $1$.
  \item $y=(x+2)^2-5$: \; vertex $(\blank{0.7cm},\,\blank{0.7cm})$ --- the parent moved
        \blank{1.6cm} $2$ and \blank{1.4cm} $5$.
  \item Fill in the rule you have used since Unit 1: a number attached \emph{outside} the function
        moves the graph \blank{2.0cm} and does exactly what it says; a number attached
        \emph{inside} moves the graph \blank{2.0cm} and does the \blank{2.0cm} of what it says.
\end{enumerate}

\vspace{4pt}
\textbf{3. One function, two costumes} (Lessons 4.3 and 4.0).
\begin{enumerate}[label=(\alph*), leftmargin=2.1em, itemsep=3pt]
  \item Combine into a single fraction: \; $2+\dfrac{3}{x-1}=\dfrac{2(\rule{1.2cm}{0.4pt})+3}{x-1}
        =$ \blank{2.0cm}
  \item For that single fraction, set the denominator to zero: vertical asymptote
        $x=$ \blank{1.0cm}.
  \item Its numerator and denominator have the \emph{same} degree, so by the degree comparison the
        horizontal asymptote is $y=$ \blank{1.0cm} (the ratio of the leading coefficients).
\end{enumerate}

\end{notesbox}

\end{document}
